% GENERATED FILE. Edit canonical lessons, then run npm run generate:books.
% canonical-source-hash: fnv1a64:226c26e8e4fbd543
% canonical-lessons: TE-C284-vadulu, TE-C284-tavvu, TE-C284-jarugu, TE-C284-cimmu, TE-C284-pisuku

\chapter{To let go, To dig, To happen, To sweep, To knead}
\label{ch:te-to-let-go-to-dig-to-happen-to-sweep-to-knead}
\hlchaptermodality{284}

\begin{chapteropening}
\textbf{By the end of this chapter:} \emph{I can say to let go, to dig, to happen, to sweep and to knead.}

Complete the last lesson of chapter 284: I can say to let go, to dig, to happen, to sweep and to knead.
\end{chapteropening}

\section[\texorpdfstring{\te{వదులు} (vadulu)}{వదులు (vadulu)}]{\te{వదులు} (vadulu) — to let go, to leave}

\label{lesson:TE-C284-vadulu}



\begin{quote}
Before the new one: say the Telugu for to oppose, then the Telugu for to argue.
\end{quote}

\subsection*{You'll want to know: \te{వదులు}}
\textbf{\te{వదులు}} — \emph{vadulu} — \textquotedblleft{}to let go, to leave\textquotedblright{}.

\begin{grammarlens}[title={What you've built}]
One more word for this chapter.
\end{grammarlens}

\subsection*{Your turn}
\begin{itemize}
\raggedright
  \item \emph{Say it:} \emph{vadulu}
  \item \emph{Say it:} \emph{vadulu}, once more
  \item \emph{Say it:} say \emph{vadulu}
  \item \emph{Recall:} say the Telugu for to oppose, then the Telugu for to argue, then say \emph{vadulu} again
\end{itemize}

\begin{tcolorbox}[breakable,colback=teal!4,colframe=teal!35!black,title={Before you move on}]
What does \te{వదులు} mean? (\textquotedblleft{}To let go, to leave\textquotedblright{}.) Say it once more.
\end{tcolorbox}
\section[\texorpdfstring{\te{తవ్వు} (tavvu)}{తవ్వు (tavvu)}]{\te{తవ్వు} (tavvu) — to dig}

\label{lesson:TE-C284-tavvu}



\begin{quote}
Before the new one: say the Telugu for to argue, then the Telugu for to let go.
\end{quote}

\subsection*{You'll want to know: \te{తవ్వు}}
\textbf{\te{తవ్వు}} — \emph{tavvu} — \textquotedblleft{}to dig\textquotedblright{}.

\begin{grammarlens}[title={What you've built}]
One more word for this chapter.
\end{grammarlens}

\subsection*{Your turn}
\begin{itemize}
\raggedright
  \item \emph{Say it:} \emph{tavvu}
  \item \emph{Say it:} \emph{tavvu}, once more
  \item \emph{Say it:} say \emph{tavvu}
  \item \emph{Recall:} say the Telugu for to argue, then the Telugu for to let go, then say \emph{tavvu} again
\end{itemize}

\begin{tcolorbox}[breakable,colback=teal!4,colframe=teal!35!black,title={Before you move on}]
What does \te{తవ్వు} mean? (\textquotedblleft{}To dig\textquotedblright{}.) Say it once more.
\end{tcolorbox}
\section[\texorpdfstring{\te{జరుగు} (jarugu)}{జరుగు (jarugu)}]{\te{జరుగు} (jarugu) — to happen}

\label{lesson:TE-C284-jarugu}



\begin{quote}
Before the new one: say the Telugu for to let go, then the Telugu for to dig.
\end{quote}

\subsection*{You'll want to know: \te{జరుగు}}
\textbf{\te{జరుగు}} — \emph{jarugu} — \textquotedblleft{}to happen\textquotedblright{}.

\begin{grammarlens}[title={What you've built}]
One more word for this chapter.
\end{grammarlens}

\subsection*{Your turn}
\begin{itemize}
\raggedright
  \item \emph{Say it:} \emph{jarugu}
  \item \emph{Say it:} \emph{jarugu}, once more
  \item \emph{Say it:} say \emph{jarugu}
  \item \emph{Recall:} say the Telugu for to let go, then the Telugu for to dig, then say \emph{jarugu} again
\end{itemize}

\begin{tcolorbox}[breakable,colback=teal!4,colframe=teal!35!black,title={Before you move on}]
What does \te{జరుగు} mean? (\textquotedblleft{}To happen\textquotedblright{}.) Say it once more.
\end{tcolorbox}
\section[\texorpdfstring{\te{చిమ్ము} (cimmu)}{చిమ్ము (cimmu)}]{\te{చిమ్ము} (cimmu) — to sweep (the floor)}

\label{lesson:TE-C284-cimmu}



\begin{quote}
Before the new one: say the Telugu for to dig, then the Telugu for to happen.
\end{quote}

\subsection*{You'll want to know: \te{చిమ్ము}}
\textbf{\te{చిమ్ము}} — \emph{cimmu} — \textquotedblleft{}to sweep (the floor)\textquotedblright{}.

\begin{grammarlens}[title={What you've built}]
One more word for this chapter.
\end{grammarlens}

\subsection*{Your turn}
\begin{itemize}
\raggedright
  \item \emph{Say it:} \emph{cimmu}
  \item \emph{Say it:} \emph{cimmu}, once more
  \item \emph{Say it:} say \emph{cimmu}
  \item \emph{Recall:} say the Telugu for to dig, then the Telugu for to happen, then say \emph{cimmu} again
\end{itemize}

\begin{tcolorbox}[breakable,colback=teal!4,colframe=teal!35!black,title={Before you move on}]
What does \te{చిమ్ము} mean? (\textquotedblleft{}To sweep (the floor)\textquotedblright{}.) Say it once more.
\end{tcolorbox}
\section[\texorpdfstring{\te{పిసుకు} (pisuku)}{పిసుకు (pisuku)}]{\te{పిసుకు} (pisuku) — to knead, to squeeze}

\label{lesson:TE-C284-pisuku}



\begin{quote}
Before the new one: say the Telugu for to happen, then the Telugu for to sweep.
\end{quote}

\subsection*{You'll want to know: \te{పిసుకు}}
\textbf{\te{పిసుకు}} — \emph{pisuku} — \textquotedblleft{}to knead, to squeeze\textquotedblright{}.

\begin{grammarlens}[title={What you've built}]
One more word for this chapter.
\end{grammarlens}

\subsection*{Your turn}
\begin{itemize}
\raggedright
  \item \emph{Say it:} \emph{pisuku}
  \item \emph{Say it:} \emph{pisuku}, once more
  \item \emph{Say it:} say \emph{pisuku}
  \item \emph{Recall:} say the Telugu for to happen, then the Telugu for to sweep, then say \emph{pisuku} again
\end{itemize}

\begin{tcolorbox}[breakable,colback=teal!4,colframe=teal!35!black,title={Before you move on}]
What does \te{పిసుకు} mean? (\textquotedblleft{}To knead, to squeeze\textquotedblright{}.) Say it once more.
\end{tcolorbox}
